% Generated draft: requires mathematical and source review.
\documentclass[11pt]{article}
\usepackage[margin=25mm]{geometry}
\usepackage{fontspec}
\setmainfont{Latin Modern Roman}
\usepackage{amsmath,amssymb,mathtools,unicode-math}
\setmathfont{Latin Modern Math}
\usepackage{enumitem,needspace,xurl,hyperref}
\setlength{\parindent}{0pt}
\setlength{\parskip}{5pt}
\setlength{\emergencystretch}{3em}
\newcommand{\N}{\mathbb{N}}
\newcommand{\Z}{\mathbb{Z}}
\newcommand{\Q}{\mathbb{Q}}
\newcommand{\R}{\mathbb{R}}
\newcommand{\C}{\mathbb{C}}
\newcommand{\F}{\mathbb{F}}
\newcommand{\A}{\mathbb{A}}
\newcommand{\PP}{\mathbb{P}}
\newcommand{\E}{\mathbb{E}}
\newcommand{\eps}{\varepsilon}
\newcommand{\dd}{\,\mathrm{d}}
\begin{document}
\section*{UnsolvedMath ID: 800321}
\textbf{Source problem code:} catalogue entry
\subsection{Definitions and mathematical statement}
The source applies differential terminology to sequences without specifying the associated generating functions or coefficient domain. The following is a conditional formulation using ordinary generating functions; these conventions require confirmation.

Fix a field \(K\) of characteristic zero and interpret differential algebraicity over \(K(t)\). The source does not specify this field or base field. For \(H(t)\in K[[t]]\), write \(H^{(j)}\) for its \(j\)-th formal derivative. Call \(H\) \emph{D-algebraic} if there exist an integer \(r\geq 0\) and a nonzero polynomial \(P\in K(t)[X_0,\ldots,X_r]\) such that
\[
P\bigl(H,H',\ldots,H^{(r)}\bigr)=0.
\]
Call \(H\) \emph{D-transcendental} if no such polynomial exists. Call a sequence \((b_n)_{n\geq 0}\) in \(K\) \emph{P-recursive} if there exist an integer \(s\geq 0\) and polynomials \(q_0,\ldots,q_s\in K[n]\), not all zero, such that
\[
\sum_{j=0}^{s}q_j(n)b_{n+j}=0
\qquad\text{for all sufficiently large integers }n.
\]

\textbf{Question under these conventions.} For every sequence \((a_n)_{n\geq 0}\) in the fixed field \(K\), set
\[
b_n=\frac{a_n}{n!},\qquad
F(t)=\sum_{n\geq 0}\frac{a_n}{n!}t^n,
\qquad
G(t)=\sum_{n\geq 0}a_nt^n.
\]
If \(F(t)\) is D-algebraic and \((b_n)_{n\geq 0}\) is not P-recursive, must \(G(t)\) be D-transcendental?

This formulation interprets the source's assertion that \(\{a_n/n!\}\) is D-algebraic as a statement about \(F\), and its proposed D-transcendence of \(\{a_n\}\) as a statement about \(G\). Neither interpretation is explicitly defined in the supplied excerpt. The choice of indexing from zero is also a convention. All series here are formal; the source supplies no convergence hypotheses.

The supplied background describes the question as an advanced generalization of Klazar's theorem and associates it with Pak--Yeliussizov. It also discusses a function \(y(t)\), quantities \(B_n\), and a factorial transformation of a P-recursive sequence. Those objects are undefined in the excerpt, and the displayed transformed sequence is truncated. Consequently, that background cannot be reconstructed as a precise additional assertion from the supplied evidence.
\subsection{Short English statement}
Under the stated generating-function interpretation, ask whether a D-algebraic exponential generating function with non-P-recursive coefficients forces the corresponding ordinary generating function to be D-transcendental. The coefficient domain and generating-function conventions remain unresolved. The supplied registry labels the problem open; this label has not been independently verified.
\subsection{Sources}
\begin{itemize}
\item Local UnsolvedMath registry, permanent ID \(800321\), exact source problem code \texttt{catalogue entry}, collection \emph{AMR Open Problem Lists}. The supplied record attributes the proposal to Pak, gives the year 2018, and labels its status \texttt{open}. Upstream revision: \texttt{d371ee917ac9647fbb06cd6fc529deec6f3dda82}.
\item Supplied source list: Pak, \emph{Complexity Problems in Enumerative Combinatorics} (2018), Problem 5, \url{https://arxiv.org/abs/1803.06636}. The record reports access on 2026-10-16 and extraction by \texttt{arxiv\_latex}; these are supplied provenance details, not an independent source visit or verification.
\item Supplied attribution/context: Pak--Yeliussizov. The excerpt mentions Klazar's theorem and the reference label [LR], but supplies neither the theorem's statement nor bibliographic details identifying [LR].
\end{itemize}
\end{document}
